\section{Постановка}
\label{sec:setup}

Частица в накопителе совершает синхротронные колебания около синхронной частицы. Если период обращения зависит от отклонения энергии нелинейно или от амплитуды бетатронных колебаний, среднее отклонение энергии за синхротронный период смещается от нуля. Это среднее --- равновесный сдвиг $\delta_\eq$ --- и находится ниже. Схема такая: сначала строится функция $\tau$, относительное опоздание частицы за оборот (разделы~\ref{sec:mag} и~\ref{sec:beta}); затем один раз решаются продольные уравнения с произвольной $\tau$ (раздел~\ref{sec:sol}); электростатическое кольцо меняет только коэффициенты $\tau$ (раздел~\ref{sec:es}).

\subsection{Переменная $\delta$}

Индекс $s$ отмечает синхронную частицу: импульс $p_s$, скорость $v_s=\beta_s c$, лоренц-фактор $\gamma_s$, энергия $E_s=\gamma_s mc^2$, длина орбиты $C_s$, период $T_s=C_s/v_s$. Функции Твисса $\beta_{x,y}$, $\alpha_{x,y}$, $\gamma_{x,y}$ всегда пишутся с индексом плоскости.

Между ускоряющими зазорами поле статическое, и сохраняется полная энергия $E=\gamma mc^2+q\Phi$. Отклонение $\delta$ определяем через неё: $p_s(1+\delta)$ есть импульс, который частица с полной энергией $E$ имела бы там, где $\Phi=0$,
\begin{equation}
\label{eq:set-delta}
E=\sqrt{p_s^2(1+\delta)^2c^2+m^2c^4}.
\end{equation}
В магнитном кольце $\Phi\equiv 0$, и $\delta=(p-p_s)/p_s$ --- обычное отклонение импульса. В электростатическом кольце локальный импульс в дефлекторе меняется вдоль траектории, а $\delta$ между зазорами постоянно и меняется только в зазоре. Поэтому именно $\delta$ --- медленная переменная продольного движения, и равновесный сдвиг есть её среднее по синхротронному колебанию:
\begin{equation}
\label{eq:set-deq}
\delta_\eq=\langle\delta\rangle.
\end{equation}

Длина замкнутой орбиты записывается рядом с коэффициентами $\alpha_0$, $\alpha_1$ (коэффициенты ряда, не производные):
\begin{equation}
\label{eq:set-C}
\frac{C}{C_s}=1+\alpha_0\delta+\alpha_1\delta^2+O(\delta^3).
\end{equation}

\subsection{Опоздание за оборот}

Относительное опоздание
\begin{equation}
\label{eq:set-tau}
\tau=\frac{T-T_s}{T_s}=\eta_0\delta+\eta_1\delta^2+\Lambda
\end{equation}
есть ряд по $\delta$ плюс вклад $\Lambda$ бетатронных колебаний, усреднённый по быстрой бетатронной фазе. Делитель --- синхронный период: за оборот фаза частицы относительно поля смещается ровно на $\omega_\rf(T-T_s)$, и в фазовое уравнение входит именно $(T-T_s)/T_s$.

\zamech{Если коэффициент скольжения определять через $(T-T_s)/T$ или через частоту обращения, $\eta=-\Delta\omega/(\omega_0\delta)$ (формула (18.20) в~\cite{Ng2006}), то во втором порядке вместо $\eta_1$ получается $\eta_1-\eta_0^2$; определение через производную, $\eta=-\omega_0^{-1}\mathrm{d}\omega/\mathrm{d}\delta$ (18.21), даёт $2(\eta_1-\eta_0^2)$. Линейный коэффициент у всех определений общий, но в фазовое уравнение входит только $\eta_1$ из~\eqref{eq:set-tau}.}

\subsection{Уравнения движения}

Фаза $\varphi$ --- отставание от синхронной частицы в радианах ВЧ-поля, она растёт при опоздании. Напряжение зазора $V\sin(\varphi_s+\varphi)$, $\omega_\rf=2\pi h/T_s$. Без ускорения $\varphi_s=0$ или $\pi$; устойчива точка
\begin{equation}
\label{eq:set-phis}
\cos\varphi_s=-\operatorname{sgn}\eta_0 .
\end{equation}
За зазор полная энергия меняется на $eV\cos\varphi_s\sin\varphi$, а $\mathrm{d}E=\beta_s^2E_s\,\mathrm{d}\delta$. Сглаживание по обороту даёт
\begin{equation}
\label{eq:set-system}
\dot\varphi=\omega_\rf\,\tau,
\qquad
\dot\delta=K\cos\varphi_s\sin\varphi,
\qquad
K=\frac{eV\omega_\rf}{2\pi h\,\beta_s^2E_s}.
\end{equation}
Точка обозначает производную по времени синхронной частицы. Линеаризация около $\varphi=\delta=0$ даёт колебания с частотой
\begin{equation}
\label{eq:set-ws}
\omega_s^2=K\omega_\rf|\eta_0|
\end{equation}
именно при выборе~\eqref{eq:set-phis}.
